\ifx\XeTeXversion\undefined\pdfoutput=1\fi
\documentclass[11pt]{article}

\usepackage[margin=1.1in]{geometry}
\ifx\XeTeXversion\undefined
  \usepackage[T1]{fontenc}
  \usepackage{lmodern}
  \usepackage[expansion=false]{microtype}
\fi
\usepackage{amsmath,amsthm,amssymb}
\usepackage{booktabs}
\usepackage{url}
\usepackage[colorlinks=true,linkcolor=black,citecolor=blue,urlcolor=blue]{hyperref}

\emergencystretch=1.5em
\clubpenalty=10000
\widowpenalty=10000

% === AUTHOR CONFIGURATION ===
\newcommand{\AuthorName}{Katsutoshi Hashimoto}
\newcommand{\AuthorAffiliation}{Independent Researcher, Japan}
\newcommand{\AuthorEmail}{ramuza0113@gmail.com}
\newcommand{\AuthorORCID}{0009-0009-8292-5085}
% ============================

\newtheorem{theorem}{Theorem}
\newtheorem{proposition}[theorem]{Proposition}
\newtheorem{lemma}[theorem]{Lemma}
\theoremstyle{remark}
\newtheorem{remark}[theorem]{Remark}

\newcommand{\F}{\mathbb{F}}
\newcommand{\Z}{\mathbb{Z}}

\hypersetup{
  pdftitle={An Improved Lower Bound for the Hypergraph Ramsey Number R(6,6;4)},
  pdfauthor={\AuthorName},
}

\title{\bfseries An Improved Lower Bound for the\\ Hypergraph Ramsey Number $R(6,6;4)$}

\author{\AuthorName\\[4pt]
\small \AuthorAffiliation\\[2pt]
\small\texttt{\AuthorEmail} \;\textperiodcentered\; \href{https://orcid.org/\AuthorORCID}{ORCID: \texttt{\AuthorORCID}}}

\date{November 2026}

\begin{document}
\maketitle

\begin{abstract}
Let $R(k,l;t)$ denote the smallest integer $n$ such that every $2$-colouring
of the $t$-subsets of an $n$-element set contains either a $k$-subset all of
whose $t$-subsets have colour~$0$ or an $l$-subset all of whose $t$-subsets
have colour~$1$. Extending projective-line constructions from
$\mathrm{PGL}(2,q)$ to the index-$2$ subgroup $\mathrm{PSL}(2,q)$ for primes
$q\equiv 3\pmod 4$, we exhibit an explicit $\mathrm{PSL}(2,619)$-invariant
$2$-colouring of the $4$-subsets of the $620$-point projective line
$\F_{619}\cup\{\infty\}$ with no monochromatic $6$-subset in either colour.
This establishes
\[
  R(6,6;4)\ge 621,
\]
improving the previous bounds $133$ listed in Radziszowski's dynamic survey,
$317$ obtained recently by \"Ozdemir, and $601$ established via
$\mathrm{PGL}(2,599)$ by Lysenst{\o}en. An independent clean-room verification
program and machine-readable witness table accompany this note.
\end{abstract}

\medskip
\noindent\textbf{Keywords:} hypergraph Ramsey numbers; constructive lower
bounds; projective special linear group; cross-ratio; stochastic local search;
computational verification.\\
\textbf{MSC 2020:} 05D10; 05C55; 05C65.

\section{Introduction}

How large can a finite set be if its $4$-element subsets are coloured with two
colours so that no $6$ points have all $\binom{6}{4}=15$ of their $4$-subsets
monochromatic? This question defines the diagonal $4$-uniform hypergraph Ramsey
number $R(6,6;4)$---more generally, $R(k,l;t)$ is the least integer $n$ such
that every $2$-colouring of $\binom{[n]}{t}$ yields a colour-$0$ $k$-subset or
a colour-$1$ $l$-subset. Since Erd\H{o}s, Hajnal, and Rado initiated the study
of hypergraph Ramsey numbers in the 1960s and McKay and Radziszowski determined
$R(4,4;3)=13$, constructive lower bounds for small $t$-uniform Ramsey numbers
have served as a central benchmark connecting finite geometry, permutation
groups, and combinatorial search~\cite{survey}.

For $t=4$, exact values remain out of reach even for $R(5,5;4)\ge 49$~\cite{ozdemir},
while $R(6,6;4)$ has undergone striking constructive progress in 2026: the
long-standing bound $R(6,6;4)\ge 133$ in Revision~\#18 of Radziszowski's
dynamic survey~\cite{survey} was raised to $317$ by \"Ozdemir~\cite{ozdemir}
and very recently to $601$ by Lysenst{\o}en~\cite{lysenstoen} using a
$\mathrm{PGL}(2,599)$-invariant colouring of the projective line
$\F_{599}\cup\{\infty\}$.

The purpose of this note is both to improve this bound and to explain two
structural mechanisms governing projective-line Ramsey constructions for
$4$-subsets. In~\cite{lysenstoen}, $\mathrm{PSL}(2,q)$-invariant colourings
were employed for the off-diagonal number $R(5,6;4)$ at $q=199$, whereas only
full $\mathrm{PGL}(2,q)$ invariance was searched for $R(6,6;4)$. First, we note
(Lemma~\ref{lem:3hom}) that whenever $q\equiv 3\pmod 4$, restricting from
$\mathrm{PGL}(2,q)$ to its index-$2$ subgroup $\mathrm{PSL}(2,q)$ preserves
transitivity on unordered $3$-subsets of the projective line---keeping the
normalized verification space at exact size $\binom{q-2}{3}$---while splitting
roughly one quarter of the cross-ratio orbits in two. Second, by Dickson's
classification of subgroups of $\mathrm{PSL}(2,q)$~\cite{dickson}, octahedral
subgroups isomorphic to $S_4$ occur in $\mathrm{PSL}(2,q)$ if and only if
$q\equiv\pm 1\pmod 8$; passing from $q\equiv 7\pmod 8$ to primes
$q\equiv 3\pmod 8$ eliminates $6$-point octahedral orbits whose $15$
$4$-subsets collapse into short symmetric clauses (Remark~\ref{rem:dickson}).
Combining these two observations with a cache-resident hybrid stochastic local
search yields an explicit witness at $q=619=8\cdot 77+3$.

\begin{theorem}\label{thm:main}
$R(6,6;4)\ge 621$.
\end{theorem}

Table~\ref{tab:history} summarizes the progression of constructive lower bounds
for $R(6,6;4)$.

\begin{table}[ht]
\centering
\small
\begin{tabular}{llccc}
\toprule
Source & Invariance group & Order $n$ & Orbit variables & Lower bound\\
\midrule
Dynamic Survey DS1~\cite{survey} & Classical / stepping-up & $132$ & -- & $133$\\
\"Ozdemir~\cite{ozdemir} & $\mathrm{PSL}(2,q)$ / projective & $316$ & -- & $317$\\
Lysenst{\o}en~\cite{lysenstoen} & $\mathrm{PGL}(2,599)$ & $600$ & $100$ & $601$\\
Present work ($F_{620}$) & $\mathrm{PSL}(2,619)$ & $620$ & $130$ & $\mathbf{621}$\\
\bottomrule
\end{tabular}
\caption{Constructive lower bounds for the $4$-uniform diagonal Ramsey number $R(6,6;4)$.}
\label{tab:history}
\end{table}

\section{Projective Line Orbits Under $\mathrm{PSL}(2,q)$ for $q\equiv 3\pmod 4$}\label{sec:psl}

Let $q$ be an odd prime and let $L_q=\F_q\cup\{\infty\}$ denote the projective
line of cardinality $q+1$. The group $\mathrm{PGL}(2,q)$ acts sharply
$3$-transitively on ordered triples of distinct points of $L_q$ by fractional
linear transformations $z\mapsto(\alpha z+\beta)/(\gamma z+\delta)$ with
$\alpha\delta-\beta\gamma\in\F_q^\times$, and $\mathrm{PSL}(2,q)$ is the
index-$2$ subgroup consisting of classes represented by matrices whose
determinant $\alpha\delta-\beta\gamma$ is a non-zero square in $\F_q$ (or
equivalently equal to~$1$).

\begin{lemma}\label{lem:3hom}
If $q\equiv 3\pmod 4$, then $\mathrm{PSL}(2,q)$ acts transitively on the
$\binom{q+1}{3}$ unordered $3$-subsets of $L_q$. Consequently, a
$\mathrm{PSL}(2,q)$-invariant $2$-colouring of $\binom{L_q}{4}$ has no
monochromatic $6$-subset if and only if none of the $\binom{q-2}{3}$
$6$-subsets containing $\{\infty,0,1\}$ is monochromatic.
\end{lemma}

\begin{proof}
Given any $3$-subset $T=\{u,v,w\}\subset L_q$, sharp $3$-transitivity of
$\mathrm{PGL}(2,q)$ yields a unique map $g\in\mathrm{PGL}(2,q)$ sending the
ordered triple $(u,v,w)$ to $(\infty,0,1)$. Let $h\in\mathrm{PGL}(2,q)$ be the
involution $z\mapsto 1-z$, represented by $\bigl(\begin{smallmatrix}-1&1\\0&1\end{smallmatrix}\bigr)$
of determinant $-1$, which swaps $0$ and $1$ and fixes $\infty$. Because
$q\equiv 3\pmod 4$, $-1$ is a quadratic non-residue in $\F_q^\times$. Thus
exactly one of $g$ and $h\circ g$ has square determinant and therefore belongs
to $\mathrm{PSL}(2,q)$, mapping $T$ setwise onto $\{\infty,0,1\}$. Since every
map in $\mathrm{PSL}(2,q)$ preserves colours of $4$-subsets, any monochromatic
$6$-subset $S\subset L_q$ is carried by an element of $\mathrm{PSL}(2,q)$
sending three points of $S$ onto $\{\infty,0,1\}$ to a monochromatic $6$-subset
containing $\{\infty,0,1\}$.
\end{proof}

To parametrize the orbits of $\mathrm{PSL}(2,q)$ on $\binom{L_q}{4}$ when
$q\equiv 3\pmod 4$, identify $L_q$ with $\{0,1,\dots,q-1,\infty\}$ ordered so
that $0<1<\dots<q-1<\infty$. For $a<b<c<d$ in $L_q$, define
\begin{equation}\label{eq:rule}
  x = \begin{cases}
    \dfrac{(a-c)(b-d)}{(a-d)(b-c)} \bmod q & \text{if } d<\infty,\\[10pt]
    \dfrac{a-c}{b-c} \bmod q & \text{if } d=\infty,
  \end{cases}
  \qquad
  \Delta = (b-a)(c-a)(c-b) \bmod q.
\end{equation}
Notice that the matrix sending $(a,b,c)\mapsto(0,1,\infty)$ or $(\infty,0,1)$
has determinant up to squares governed by the Legendre symbol of $\Delta$:
specifically, the unique fractional linear transformation sending
$(a,b,c,d)\mapsto(0,1,x,\infty)$ has determinant congruent modulo squares to
$-\Delta$. Under the alternating subgroup $A_3=\{\mathrm{id},\,z\mapsto 1/(1-z),\,z\mapsto(z-1)/z\}\le\mathrm{PSL}(2,q)$
stabilizing $\{\infty,0,1\}$, each $x\in\F_q\setminus\{0,1\}$ generates a
$3$-element orbit (or $1$-element orbit if $x^2-x+1\equiv 0\pmod q$)
\[
  C_3(x)=\Bigl\{\,x,\;\frac{1}{1-x},\;\frac{x-1}{x}\,\Bigr\}\subset\F_q\setminus\{0,1\}.
\]
Meanwhile, the Klein four-group stabilizing $\{\infty,0,1,x\}$ in
$\mathrm{PGL}(2,q)$ is represented by matrices of determinants
$1$, $-x$, $x-1$, and $-x(x-1)$. All four lie in $\mathrm{PSL}(2,q)$ if and
only if both $-x$ and $x-1$ are quadratic residues modulo $q$; otherwise the
remaining two elements of $\mathrm{Stab}_{\mathrm{PGL}(2,q)}(\{\infty,0,1,x\})$
have non-square determinant, fusing $C_3(x)$ with $C_3(1/x)=\{1/x,\,1-x,\,x/(x-1)\}$
under $\mathrm{PSL}(2,q)$. Hence the $\mathrm{PSL}(2,q)$-orbits on
$\binom{L_q}{4}$ correspond bijectively to the \emph{narrow classes}
\[
  N(x)=\begin{cases}
    C_3(x) & \text{if } \left(\frac{-x}{q}\right)=+1 \text{ and } \left(\frac{x-1}{q}\right)=+1,\\[6pt]
    C_3(x)\cup C_3(1/x) & \text{otherwise,}
  \end{cases}
\]
and for any table $\sigma\colon\F_q\setminus\{0,1\}\to\{0,1\}$ constant on
narrow classes, assigning to $\{a,b,c,d\}$ ($a<b<c<d$) the colour
\begin{equation}\label{eq:col}
  \chi_\sigma(\{a,b,c,d\})=\begin{cases}
    \sigma(x) & \text{if $\Delta$ is a square modulo $q$,}\\
    \sigma(x^{-1}\bmod q) & \text{otherwise}
  \end{cases}
\end{equation}
defines a $\mathrm{PSL}(2,q)$-invariant $2$-colouring of $\binom{L_q}{4}$
(matching the convention used for $F_{200}$ in~\cite{lysenstoen}).

\begin{remark}[Subgroup obstructions and $q\equiv 3\pmod 8$]\label{rem:dickson}
By Dickson's classification of subgroups of $\mathrm{PSL}(2,q)$~\cite{dickson}, the symmetric
group $S_4$ embeds into $\mathrm{PSL}(2,q)$ if and only if $q\equiv\pm 1\pmod 8$.
When $q\equiv 7\pmod 8$ (such as $q=599,607,631,647$), octahedral $S_4$-orbits
of size~$6$ in $L_q$ partition their $15$ $4$-subsets into very few
$\mathrm{PSL}(2,q)$-orbits, generating dozens of short symmetric clauses of
width $\le 6$. Choosing primes $q\equiv 3\pmod 8$ (which automatically satisfy
$q\equiv 3\pmod 4$ while excluding $S_4\le\mathrm{PSL}(2,q)$), such as
$q=619=8\cdot 77+3$, eliminates these octahedral short clauses while retaining
roughly $1.25\times$ as many orbit variables as $\mathrm{PGL}(2,q)$.
\end{remark}

\section{The Witness $F_{620}$ for $q=619$}\label{sec:witness}

Let $q=619$. Since $619$ is prime and $619\equiv 3\pmod 8$, $\F_{619}\setminus\{0,1\}$
partitions into $130$ narrow classes under $\mathrm{PSL}(2,619)$ (comprising
$53$ classes of size~$3$ and $77$ classes of size~$6$, whereas
$\mathrm{PGL}(2,619)$ has $104$ classes).

We define the colouring $F_{620}$ on $L_{619}$ ($620$ points,
$\binom{620}{4}=6\,083\,590\,570$ $4$-subsets) via~\eqref{eq:rule}--\eqref{eq:col}
by setting $\sigma=0$ on the following $62$ narrow classes (each identified by
its smallest representative $d\in\{2,\dots,618\}$) and $\sigma=1$ on the
remaining $68$ narrow classes:
\begin{quote}
\small\raggedright
$d = 2$,\allowbreak\ $4$,\allowbreak\ $5$,\allowbreak\ $6$,\allowbreak\ $7$,\allowbreak\ $9$,\allowbreak\ $10$,\allowbreak\ $14$,\allowbreak\ $15$,\allowbreak\ $16$,\allowbreak\ $18$,\allowbreak\ $26$,\allowbreak\ $27$,\allowbreak\ $28$,\allowbreak\ $30$,\allowbreak\ $32$,\allowbreak\ $38$,\allowbreak\ $39$,\allowbreak\ $41$,\allowbreak\ $43$,\allowbreak\ $44$,\allowbreak\ $45$,\allowbreak\ $47$,\allowbreak\ $48$,\allowbreak\ $50$,\allowbreak\ $51$,\allowbreak\ $52$,\allowbreak\ $57$,\allowbreak\ $61$,\allowbreak\ $62$,\allowbreak\ $66$,\allowbreak\ $67$,\allowbreak\ $68$,\allowbreak\ $69$,\allowbreak\ $70$,\allowbreak\ $71$,\allowbreak\ $72$,\allowbreak\ $74$,\allowbreak\ $78$,\allowbreak\ $81$,\allowbreak\ $82$,\allowbreak\ $84$,\allowbreak\ $96$,\allowbreak\ $102$,\allowbreak\ $108$,\allowbreak\ $109$,\allowbreak\ $110$,\allowbreak\ $118$,\allowbreak\ $122$,\allowbreak\ $126$,\allowbreak\ $137$,\allowbreak\ $158$,\allowbreak\ $167$,\allowbreak\ $174$,\allowbreak\ $175$,\allowbreak\ $208$,\allowbreak\ $213$,\allowbreak\ $215$,\allowbreak\ $230$,\allowbreak\ $249$,\allowbreak\ $250$,\allowbreak\ $321$.
\par
\end{quote}
Across $x\in\{2,\dots,618\}$, $\sigma(x)$ takes value $0$ at $303$ arguments
and value $1$ at $314$ arguments. Extending $\sigma$ by $\sigma(0)=\sigma(1)=0$
to a $619$-character ASCII string \texttt{table} of \texttt{0}s and \texttt{1}s,
its SHA-256 digest is
\[
  \texttt{6658e8cfc640c17e14b425f5b0a19bd03017ed1aba7a035752df5550199467c6}.
\]

\begin{proposition}\label{prop:valid}
$F_{620}$ is a $(6,6;4)$-colouring of order $620$; that is, none of the
$\binom{620}{6}=79\,532\,803\,405\,740$ $6$-subsets of $L_{619}$ has all $15$
of its $4$-subsets in colour~$0$ or all $15$ in colour~$1$.
\end{proposition}

\begin{proof}[Computational verification]
Established by the independent exhaustive verification described in
Section~\ref{sec:verif}.
\end{proof}

Theorem~\ref{thm:main} follows immediately from Proposition~\ref{prop:valid}.

\section{Search Method and Independent Verification}\label{sec:verif}

\paragraph{Stochastic local search.}
For $q=619$, deduplicating the orbit bitmasks of the $\binom{617}{3}=38\,957\,380$
normalized $6$-subsets $\{\infty,0,1,x,y,z\}$ ($2\le x<y<z<619$) produces
$620\,672$ distinct subsets of the $130$ narrow-class variables (yielding a
monotone/anti-monotone $k$-SAT instance of $1\,241\,344$ clauses of width up
to~$15$). To solve this instance efficiently without exceeding CPU cache
bandwidth, each clause's live state (number of $1$-literals and bitwise XOR of
their indices) was packed into $16$ bits, allowing incremental $O(1)$ retrieval
of unique critical variables upon transitions into or out of $1$-hot and
$(k-1)$-hot states. A portfolio combining exponential $k$-SAT
probSAT~\cite{probsat} ($f(\mathrm{break})=c_b^{-\mathrm{break}}$ with a
zero-break freebie rule) and WalkSAT/SKC with short tabu tenure located
$F_{620}$ in $69\,993$ flips ($8.95$ seconds).

\paragraph{Clean-room verification.}
To eliminate any reliance on solver data structures or group-theoretic
shortcuts, $F_{620}$ was verified by two separate programs sharing no code with
the solver (\texttt{verify\_r664\_psl.rs} and \texttt{verify\_r664\_psl.py}):
\begin{enumerate}
\item \emph{Direct $\mathrm{SL}(2,619)$ well-definedness and $3$-homogeneity.}
  Rather than invoking Lemma~\ref{lem:3hom} abstractly, \texttt{verify\_r664\_psl.rs}
  works directly with $2\times 2$ matrices $M=\bigl(\begin{smallmatrix}a&b\\c&d\end{smallmatrix}\bigr)$
  satisfying $ad-bc\equiv 1\pmod{619}$. It first confirms that on every tested
  $4$-subset, all $12$ ordered lifts in $\mathrm{SL}(2,619)$ mapping three
  points to $(\infty,0,1)$ produce identical values of $\sigma(\lambda)$. Next,
  it iterates exhaustively over all $\binom{620}{3}=39\,529\,340$ unordered
  $3$-subsets $\{u,v,w\}\subset L_{619}$ and explicitly constructs and checks a
  matrix $M\in\mathrm{SL}(2,619)$ with $\det(M)\equiv 1\pmod{619}$ and
  $M(\{u,v,w\})=\{\infty,0,1\}$. In addition, $1\,000\,000$ random pairs
  $(E,M)\in\binom{L_{619}}{4}\times\mathrm{SL}(2,619)$ were tested directly to
  confirm $\chi_\sigma(M(E))=\chi_\sigma(E)$.
\item \emph{Exhaustive $6$-subset verification.}
  The verifier examines every one of the $\binom{617}{3}=38\,957\,380$
  $6$-subsets containing $\{\infty,0,1\}$. Every such $6$-subset has between
  $1$ and $14$ of its $15$ $4$-subsets in colour~$1$, leaving zero
  monochromatic $6$-subsets in colour~$0$ and zero in colour~$1$. Combined with
  Step~1, this exhaustively covers all $\binom{620}{6}$ $6$-subsets of
  $L_{619}$. As an additional symmetry-free sanity check, $500\,000$ arbitrary
  $6$-subsets sampled uniformly from $L_{619}$ were evaluated directly from raw
  $4$-subsets without normalization, again with zero violations.
\item \emph{Cross-check with independent rule reconstructor.}
  The Python verifier \path{verify_r664_psl.py} implements~\eqref{eq:rule}--\eqref{eq:col}
  using only the Python standard library, verifies the SHA-256 digest and narrow
  class constancy of \path{R664_620_points_rule.json}, and confirms exact
  agreement with the narrow-class convention of~\cite{lysenstoen}.
\end{enumerate}
The entire Rust verification suite completes in under $8$ seconds on a desktop
CPU.

\subsection*{Data Availability}
The machine-readable witness file \path{R664_620_points_rule.json}, the
clean-room verifiers \path{verify_r664_psl.rs} and \path{verify_r664_psl.py},
and the local-search solver accompany this note as ancillary files and are
permanently archived at Zenodo (\url{https://doi.org/10.5281/zenodo.24462207})
and on GitHub (\url{https://github.com/ramuza0113-sketch/ramsey-r664-621}).

\subsection*{Acknowledgments}
The author thanks Stanis{\l}aw P.\ Radziszowski for maintaining the dynamic
survey on small Ramsey numbers and Christian Lysenst{\o}en whose recent
preprint~\cite{lysenstoen} motivated this $\mathrm{PSL}(2,q)$ investigation for
$R(6,6;4)$. In accordance with journal policy on artificial intelligence tools,
the author discloses that an AI assistant (Google DeepMind Antigravity) was
used to assist in implementing the Rust local-search and clean-room
verification programs, checking algebraic details of Lemma~\ref{lem:3hom} and
Remark~\ref{rem:dickson}, and drafting and polishing the English exposition
and \LaTeX{} formatting. Every mathematical argument, citation, and
computational output was checked against independent verification, and the
witness $F_{620}$ is deterministically verifiable in seconds from the
standalone ancillary code; the author takes full responsibility for all content
of this manuscript.

\begin{thebibliography}{9}
\small

\bibitem{probsat}
A.~Balint and U.~Sch\"oning,
\emph{Choosing probability distributions for stochastic local search and the
role of make versus break},
in: Theory and Applications of Satisfiability Testing -- SAT 2012,
Lecture Notes in Comput.\ Sci.\ 7317, Springer, 2012, 16--29.

\bibitem{dickson}
L.~E. Dickson,
\emph{Linear Groups: With an Exposition of the Galois Field Theory},
B.~G. Teubner, Leipzig, 1901 (reprinted by Dover, New York, 1958).

\bibitem{lysenstoen}
C.~Lysenst{\o}en,
\emph{New lower bounds for six small hypergraph Ramsey numbers},
preprint, November 2026, arXiv:2611.25143v1.

\bibitem{ozdemir}
K.~\"Ozdemir,
\emph{Lower bounds for the hypergraph Ramsey numbers $R(5,5;4)$ and $R(5,6;4)$
from $\mathrm{PSL}(2,q)$-invariant colourings},
data deposit, Zenodo, 27 October 2026,
\url{https://doi.org/10.5281/zenodo.23689461}.

\bibitem{survey}
S.~P. Radziszowski,
\emph{Small Ramsey numbers},
Electron.\ J.\ Combin., Dynamic Survey DS1, revision~\#18 (24 April 2026),
\url{https://doi.org/10.37236/21}.

\end{thebibliography}

\end{document}
